% !TEX program = xelatex
% ---------------------------------------------------------------------------
% 這份 .tex 是 resume pipeline 從 Notion 自動產生的，可以直接在 Overleaf 改。
%
% 1. 編譯器要用 XeLaTeX（不是 pdfLaTeX）—— 第一行的 magic comment 通常會自動選好；
%    沒有的話到 Menu → Compiler 改。xeCJK 與 fontspec 都需要它。
% 2. 字型是編譯時才決定的：找不到 Noto 就會自動退到 TeX Live 內建字型。
%    若仍因缺字型編不過，只要改下面那幾行 fallback 的字型名字就行。
% 3. 在這裡的手改只存在於 Overleaf。下一次 pipeline 產出會覆蓋這個檔案 ——
%    要永久生效請改 repo 裡的 templates/resume.tex.j2。
% ---------------------------------------------------------------------------
\documentclass[11pt,a4paper]{article}

\usepackage{fontspec}
\usepackage{xeCJK}

\IfFontExistsTF{Noto Serif}
  {\setmainfont{Noto Serif}}
  {\setmainfont{TeX Gyre Termes}}
\IfFontExistsTF{Noto Sans}
  {\setsansfont{Noto Sans}}
  {\setsansfont{TeX Gyre Heros}}
\IfFontExistsTF{Noto Serif CJK TC}
  {\setCJKmainfont{Noto Serif CJK TC}}
  {\IfFontExistsTF{uming.ttc}
    {\setCJKmainfont{uming.ttc}}
    {\setCJKmainfont{FandolSong-Regular.otf}}}
\IfFontExistsTF{Noto Sans CJK TC}
  {\setCJKsansfont{Noto Sans CJK TC}}
  {\IfFontExistsTF{ukai.ttc}
    {\setCJKsansfont{ukai.ttc}}
    {\setCJKsansfont{FandolHei-Regular.otf}}}

\IfFontExistsTF{Noto Sans CJK TC}
  {%
    \newfontfamily\symbolfont{Noto Sans CJK TC}
    \newcommand{\bulletmark}{{\symbolfont ▪}}
    \newcommand{\arrowright}{{\symbolfont →}}
    \newcommand{\arrowboth}{{\symbolfont ↔}}
  }
  {%
    \newcommand{\bulletmark}{\rule[0.05ex]{0.32em}{0.32em}}
    \newcommand{\arrowright}{\ensuremath{\rightarrow}}
    \newcommand{\arrowboth}{\ensuremath{\leftrightarrow}}
  }

\xeCJKDeclareCharClass{HalfLeft}{`—,`–}

\XeTeXlinebreaklocale "zh"
\XeTeXlinebreakskip = 0pt plus 1pt minus 0.1pt

\usepackage[
  a4paper,
  top=12mm,
  bottom=12mm,
  left=14mm,
  right=14mm
]{geometry}
\usepackage{titlesec}
\usepackage{enumitem}
\usepackage{xcolor}
\usepackage[hidelinks]{hyperref}
\usepackage{tabularx}
\usepackage{ragged2e}

\definecolor{ink}{HTML}{111111}
\definecolor{inkmuted}{HTML}{444444}
\color{ink}

\RaggedRight
\setlength{\parindent}{0pt}
\pagestyle{empty}

\titleformat{\section}
  {\normalfont\sffamily\bfseries\fontsize{10.5pt}{13pt}\selectfont}
  {}{0em}{}[{\vspace{-1.1mm}\rule{\linewidth}{0.5pt}}]
\titlespacing*{\section}{0pt}{4.2mm}{1.1mm}

\setlist[itemize]{
  leftmargin=3.4mm,
  itemsep=0.7mm,
  parsep=0pt,
  topsep=1mm,
  partopsep=0pt,
  label={\footnotesize\raisebox{0.15ex}{\bulletmark}}
}

\newcommand{\entryhead}[2]{%
  \begin{tabularx}{\linewidth}{@{}Xr@{}}
    \textbf{\fontsize{10.5pt}{13pt}\selectfont #1} &
    {\color{inkmuted}\fontsize{9.5pt}{12pt}\selectfont #2}
  \end{tabularx}%
}

\newcommand{\entrysub}[2]{%
  \begin{tabularx}{\linewidth}{@{}Xr@{}}
    {\color{inkmuted}\fontsize{9.5pt}{12pt}\selectfont #1} &
    {\color{inkmuted}\fontsize{9.5pt}{12pt}\selectfont #2}
  \end{tabularx}%
}

\fontsize{9.5pt}{14.3pt}\selectfont

\begin{document}

\begin{center}
  {\sffamily\bfseries\fontsize{17pt}{20pt}\selectfont
   吳雨諠}\\[1.2mm]
  {\color{inkmuted}\fontsize{9pt}{11pt}\selectfont
   applegirlmabel@gmail.com}
\end{center}

\section{學歷}
\entryhead{國立陽明交通大學}{2028（預計）}
\entrysub{語言學碩士}{}
\begin{itemize}
  \item 預計修習課程：音韻學、計算語言學
\end{itemize}
\vspace{2.6mm}\entryhead{國立臺灣大學}{2026}
\entrysub{外國語文學系學士}{}
\begin{itemize}
  \item 參與臺大計算語言學學程
  \item 修習課程：自然語言處理、計算語言學、音韻學、語意學
\end{itemize}

\section{工作經歷}
\entryhead{國泰金控 DDT AI}{台灣台北}
\entrysub{ML/AI 工程師實習生}{2026 年 2 月 ～ 至今}
\begin{itemize}
  \item 開發針對內部業務流程優化的 LLM agent，並設計與部署其背後的雲端架構 。負責 AI 應用層與其運行的生產環境基礎架構。
\end{itemize}
\vspace{2.6mm}\entryhead{台灣人工智慧學校（AIA）}{台灣新北市}
\entrysub{AI 開發與技術支援暑期實習生}{2025 年 6 月 ～ 2025 年 8 月}
\begin{itemize}
  \item 開發與支援並行的角色：開發AI 原生應用，同時擔任企業課程現場的第一線技術支援。
\end{itemize}

\section{專案成果}
\entryhead{Voice-Notion — Slack 語音轉 Notion 機器人雲原生架構}{國泰金控 DDT AI}
\begin{itemize}
  \item 開發 Slack bot，將 mp4 口頭工作報告轉換為可確認的 Notion 頁面與任務，並自動產出每日 standup 摘要。
  \item 將單一實例的 Next.js 服務重新架構為雲原生 ECS Fargate 部署，搭配 DynamoDB、SQS 與 EventBridge，實現水平擴展與零資料遺失的滾動更新。
  \item 以維運成本與 private subnet 條件為評估標準，比較 Lambda + API Gateway、自管 EC2、App Runner 與 ECS Fargate。
\end{itemize}
\vspace{2.6mm}\entryhead{AI News Agent — 三情境路由聊天機器人}{國泰金控 DDT AI}
\begin{itemize}
  \item 以四層架構（介面 / 應用 / 資料 / 外部服務）打造內部 AI 研究報告查詢助理。
  \item 設計三情境分類系統（LLM 意圖偵測 + Python 狀態機），對明確、模糊與超出範圍的提問進行智慧路由。
  \item 建置完整 ETL 流程（Google Sheets \arrowright{} Bedrock Embedding \arrowright{} OpenSearch），並依主題、來源、日期與訂閱層級做多維度篩選。
\end{itemize}
\vspace{2.6mm}\entryhead{BU Agent — 雙模式 BRD 協作草擬助理（網頁版）}{國泰金控 DDT AI}
\begin{itemize}
  \item 以 LangGraph + PostgresSaver 打造內部網頁應用，採分屏佈局（對話 \arrowboth{} 模式感知工作區）
  \item BU Agent 作為業務單位的單一入口（兩種模式）草擬 v1.0 BRD，再將結構化交付物交棒給下游 BA Agent。
\end{itemize}

\section{論文與演講}
\entryhead{Examining LLM's Ability to Judge Presuppositions: A Comparative Study of Zero-Shot and Self-Discover-NLI Prompting}{ROCLING 2025}
\begin{itemize}
  \item 開發並發表「Self-Discover-NLI」prompt engineering 框架，將 Self-Discover 三步驟流程（SELECT \arrowright{} ADAPT \arrowright{} APPLY）重新設計以應用於自然語言推論（NLI）任務。
  \item 於 IMPPRES 資料集進行系統性實驗（每種觸發詞 95 個任務），評估 GPT-4o 對分裂句、條件句、情態詞等語言觸發詞的處理能力。
  \item 整體準確率提升約 10\%；並透過深入誤差分析指出模型的系統性弱點。
\end{itemize}

\section{專業技能}
\begin{itemize}[label={},leftmargin=0pt,itemsep=0.7mm]
  \item \textbf{語言:}\ 中文（母語）、英文（流利）
  \item \textbf{程式語言:}\ Python
  \item \textbf{框架:}\ LangGraph、LangChain
  \item \textbf{工具:}\ AWS CLI、Docker
  \item \textbf{證照:}\ 多益（聽讀 985、口說 180、寫作 170）、AWS SAA（Solutions Architect Associate）
\end{itemize}

\end{document}